\documentclass[11pt,a4paper]{article}
\usepackage{fontspec}
\usepackage{polyglossia}
\setdefaultlanguage{russian}
\setmainfont{PT Serif}
\setsansfont{PT Sans}
\setmonofont{PT Mono}
\newfontfamily\cyrillicfont{PT Serif}
\newfontfamily\cyrillicfonttt{PT Mono}
\usepackage[margin=25mm]{geometry}
\usepackage{tabularx,longtable,fancyvrb}
\usepackage[hidelinks]{hyperref}
\setlength{\parindent}{1.25em}
\setlength{\emergencystretch}{3em}
\setlength{\parskip}{0.3em}
\renewcommand{\arraystretch}{1.25}
\hypersetup{pdftitle={ИЫ.00110-01 28 01},pdfsubject={ВЫЧИСЛИТЕЛЬНЫЙ КОМПЛЕКС "ЭЛЬБРУС 1-КБ" ОБЩЕЕ СИСТЕМНОЕ ПРОГРАММНОЕ ОБЕСПЕЧЕНИЕ ТЕХНИЧЕСКИЕ УСЛОВИЯ}}
\begin{document}
\begin{titlepage}\small
\noindent\begin{minipage}[t]{0.45\linewidth}Согласовано\\Представитель заказчика №\,2272\\В.~И.~Пономаренко\end{minipage}\hfill\begin{minipage}[t]{0.45\linewidth}\raggedleft Утверждаю\\Руководитель предприятия\\Г.~Г.~Рябов\end{minipage}
\vspace*{\fill}\begin{center}
{\large Вычислительный комплекс «Эльбрус 1-КБ»}\\[0.6em]
{\large Общее системное программное обеспечение}\\[0.6em]
{\large Технические условия}\\[0.6em]
\vspace{1em}{\large\bfseries ИЫ.00110-01 28 01}\\[0.4em]
Листов 23\end{center}\vspace*{\fill}
\begin{flushright}\begin{tabular}{r}Представители предприятия-разработчика\\Начальник отдела\\В.~Ф.~Тюрин\\Исполнитель\\С.~Э.~Головчик\\Нормоконтролер\\Т.~П.~Рудаченко\end{tabular}\end{flushright}
\begin{center}1988\end{center}\end{titlepage}
\section*{Аннотация}\addcontentsline{toc}{section}{Аннотация}
Настоящие технические условия (ТУ) распространяются на общее системное программное обеспечение (ОСПО) ИЫ.00110-01, предназначенное для обеспечения прохождения задач различного назначения на вычислительном комплексе (ВК) «Эльбрус 1-КБ»; решения задач в пакетном, диалоговом и диалогово-пакетном режимах, используя при этом трансляторы с автокодов (автокод ИТМ, БЕМШ, ассемблер, MADLEN) и языков высокого уровня (Алгол, Фортран, Паскаль). Общий объем готовых программ ОСПО составляет около 600 тысяч машинных слов.

Общее системное программное обеспечение должно обеспечивать круглосуточную эксплуатацию ВК «Эльбрус 1-КБ» в нормальных климатических условиях в соответствии с требованиями ГОСТ 16325-76.

ОСПО поставляется пользователям на магнитных лентах (МЛ).

Пример записи обозначения программы при заказе:

«ИЫ.00110-01. Вычислительный комплекс «Эльбрус 1-КБ».

Общее системное программное обеспечение».

Пример записи обозначения программы в документации другой продукции:

«Вычислительный комплекс «Эльбрус 1-КБ»,

Общее системное программное обеспечение. ИЫ.00110-01».

Настоящие ТУ предназначены для работников предприятия-изготовителя, отдела технического контроля (ОТК), представителей заказчика и определяет необходимые технические требования по изготовлению, испытаниям, приемке и поставке ОСПО.

\tableofcontents\clearpage
\section*{1.~Технические требования}\addcontentsline{toc}{section}{1.~Технические требования}
\subsection*{1.1.~Общие требования}\addcontentsline{toc}{subsection}{1.1.~Общие требования}
\par\textbf{1.1.1.} Общее системное программное обеспечение ВК «Эльбрус 1-КБ» должно соответствовать требованиям настоящих ТУ и изготавливаться в соответствии с технологической документацией, утвержденной в установленном на предприятии-изготовителе ОСПО порядке.\par
\par\textbf{1.1.2.} Для изготовления ОСПО должны использоваться следующие типы магнитных носителей (МН):\par
Магнитная лента типа В4502-12 ГОСТ 20958-80. По согласованию с представителем заказчика могут использоваться и другие типы магнитных носителей, аналогичные вышеприведенным по техническим характеристикам.

\par\textbf{1.1.3.} Магнитные носители должны удовлетворять требованиям действующих на них стандартов и ТУ и должны быть приняты в установленном на предприятиях-изготовителях магнитных носителей порядке, качество магнитных носителей должно быть подтверждено соответствующими клеймами, паспортами, сертификатами качества предприятий-изготовителей.\par
Магнитные носители перед записью на них ОСПО должны пройти входной контроль в установленном порядке на предприятии-изготовителе ОСПО, согласованном с представителем заказчика.

\par\textbf{1.1.4.} Срок хранения магнитных носителей с момента их изготовления до момента записи информации на них не должен превышать 18 месяцев.\par
\par\textbf{1.1.5.} Технологический процесс изготовления ОСПО должен соответствовать согласованной и утвержденной в установленном порядке инструкции предприятия-изготовителя и обеспечивать соответствие параметров магнитных носителей с записанными на них программами и документами на ОСПО действующим нормативно-техническим документам.\par
\par\textbf{1.1.6.} Копии, дубликаты и подлинники ОСПО, расположенные на МН, должны быть идентичны.\par
\par\textbf{1.1.7.} Изготовленный экземпляр ОСПО должен быть подвергнут контролю на правильность записи на магнитные носители путем считывания со сверкой контрольных сумм (КС).\par
\par\textbf{1.1.8.} Копии ОСПО должны сниматься только с дубликата.\par
Дубликаты ОСПО должны проверяться на сохранность контрольных сумм не реже одного раза в год, о чем должны иметься соответсвующие отметки в формулярах дубликатов ОСПО.

\par\textbf{1.1.9.} Эксплуатационные программные документы на ОСПО, изготовленные на бумажном носителе, не должны иметь дефектов печати после их изготовления и должны соответствовать ведомости эксплуатационных документов.\par
\subsection*{1.2.~Комплектность}\addcontentsline{toc}{subsection}{1.2.~Комплектность}
Комплектность поставки ОСПО ВК «Эльбрус 1-КБ» должна соответствовать составу, указанному в табл. 1. ОСПО и программная документация поставляются на четырех магнитных лентах МЛ1-МЛ4, МЛ4 не входит в комплект эксплуатационных документов.

\begin{center}Состав магнитных лент:\end{center}
МЛ1 --- Программа переписи с МЛ в ОП и ГП ОС Диспак;

МЛ2 --- Системы программирования режима совместимости;

МЛ3 --- Система программирования,

Эксплуатационная документация;

МЛ4 --- Тексты программ. Данная МЛ содержит основные части документов: «Операционная система Диспак. Текст программы» ИЫ.00110-01 12 01 (файлы 1-2), «Система программирования. Текст программы» ИЫ.00110-01 12 02 файлы 3-8), «Система программирования. ассемблер. Текст программы» ИЫ.00110-01 12 03 (файлы 9-14).

Таблица 1

\begin{center}\small\begin{tabularx}{\linewidth}{|>{\raggedright\arraybackslash\hsize=1.443\hsize}X|>{\raggedright\arraybackslash\hsize=1.574\hsize}X|>{\raggedright\arraybackslash\hsize=0.393\hsize}X|>{\raggedright\arraybackslash\hsize=0.590\hsize}X|}\hline
\small\bfseries\hspace{0pt}Наименование составной части, документа & \small\bfseries\hspace{0pt}Обозначение & \small\bfseries\hspace{0pt}Количество & \small\bfseries\hspace{0pt}Примечание \\ \hline
1. Комплект программных документов согласно ведомости эксплуатационных документов & ИЫ.00110-01 20 01 & 1 & 1 \\ \hline
2. Комплект программных документов и тексты программ ОСПО на магнитных носителях с наклеенными этикетками &  &  & На магнитной ленте \\ \hline
3. Упаковка & ПЭ4.170.149-17 & 1 &  \\ \hline
\end{tabularx}\end{center}

Примечание: 1. Поставляется на бумажном носителе

\subsection*{1.3.~Маркировка}\addcontentsline{toc}{subsection}{1.3.~Маркировка}
\par\textbf{1.3.1.} Маркировка МН с ОСПО должна производиться согласно документации на магнитные ленты в соответствии с инструкцией предприятия-изготовителя ОСПО, согласованной с представителем заказчика.\par
Маркировка на магнитные ленты должна обеспечивать однозначную идентификацию класса ОСПО (копия, дубликат, подлинник), а также включать обозначение спецификации ОСПО, порядковый учетный номер магнитного носителя.

\par\textbf{1.3.2.} ОСПО, расположенное на МН в штатной (транспортной) упаковке, должно иметь клейма и пломбы в установленных местах, маркировку, определяющую содержимое упаковки, и соответствующие записи в формуляре ОСПО.\par
Вид и расположение маркировки на упаковке устанавливаются в технической документации на упаковку заводом-изготовителем упаковки и представителем заказчика.

\par\textbf{1.3.3.} Маркировка ОСПО, расположенного на МЛ, включая маркировку упаковки, должна сохраняться после нахождения в условиях, предусмотренных технической документацией на упаковку и должна соответствовать документации на тару.\par
\subsection*{1.4.~Упаковка}\addcontentsline{toc}{subsection}{1.4.~Упаковка}
\par\textbf{1.4.1.} Упаковка ОСПО должна производиться в соответствии с инструкцией предприятия-изготовителя, согласованной и утвержденной в установленном порядке.\par
\par\textbf{1.4.2.} Упаковка ОСПО должна обеспечивать сохранность эксплуатационной документации и МН при хранении и соответствовать требованиям ГОСТ 20958-80.\par
\section*{2.~Правила приемки}\addcontentsline{toc}{section}{2.~Правила приемки}
\subsection*{2.1.~Общие положения}\addcontentsline{toc}{subsection}{2.1.~Общие положения}
\par\textbf{2.1.1.} Испытания ОСПО на соответствие требованиям настоящих ТУ проводит представитель заказчика с участием ОТК предприятия-изготовителя на оборудовании, на которое имеется документация, подтверждающая его готовность к испытаниям.\par
\par\textbf{2.1.2.} Предприятие-изготовитель обязано обеспечить представителя заказчика комплектом учтенной документации на ОСПО.\par
\par\textbf{2.1.3.} ОСПО, предъявляемое представителю заказчика, должно быть принято в объеме настоящих ТУ ОТК предприятия-изготовителя и иметь клейма на МН или этикетке МН, на упаковке и в документации.\par
\par\textbf{2.1.4.} Для проверки соответствия ОСПО требованиям настоящих ТУ устанавливаются приемо-сдаточные испытания.\par
\par\textbf{2.1.5.} Испытания ОСПО проводятся в нормальных климатических условиях:\par
Температура окружающего воздуха (25±10)~°C;

Относительная влажность окружающего воздуха (65±15) \%;

Атмосферное давление (1000±40) гПа.

\par\textbf{2.1.6.} Испытания МН с ОСПО проводятся на устройствах, входящих в состав ВК «Эльбрус 1-КБ».\par
\par\textbf{2.1.7.} Формы документов, действующие на предприятии-изготовителе и применяемые во всех видах испытаний ОСПО, должны быть согласованы с представителем заказчика.\par
\subsection*{2.2.~Приемо-сдаточные испытания}\addcontentsline{toc}{subsection}{2.2.~Приемо-сдаточные испытания}
\par\textbf{2.2.1.} Приемо-сдаточные испытания проводят для проверки каждого изготовленного экземпляра ОСПО на соответствие требованиям настоящих ТУ.\par
\par\textbf{2.2.2.} Состав и последовательность приемо-сдаточных испытаний приведены в табл. 2.\par
Таблица 2

\begin{center}\small
\begin{BVerbatim}
────────────────────────────────────────────────────────────
                                     │   Номера пунктов
         наименования испытаний      │─────────────────────
                                     │технических│ методов
              и проверок             │           │
                                     │требований │испытаний
 ────────────────────────────────────│───────────│─────────
 1. Проверка комплектности           │   1.2.    │  3.2.
 ────────────────────────────────────│───────────│─────────
 2. Проверка примененного типа       │   1.1.2.  │  3.1.1.
                                     │           │
    магнитного носителя              │           │
 ────────────────────────────────────│───────────│─────────
 3. Проверка соответствия магнитного │   1.1.3.  │  3.1.2.
                                     │           │
    носителя данных требованиям      │           │
                                     │           │
    стандартов и ТУ                  │           │
 ────────────────────────────────────│───────────│─────────
 4. Проверка срока хранения магнит-  │   1.1.4.  │  3.1.3.
                                     │           │
    ного носителя                    │           │
 ────────────────────────────────────│───────────│─────────
 5. Проверка технологического про-   │   1.1.5.  │  3.1.4.
                                     │           │
    цесса изготовления ОСПО          │           │
 ────────────────────────────────────│───────────│─────────
 6. Проверка контрольных сумм        │   1.1.7.  │  3.1.5.
 ────────────────────────────────────│───────────│─────────
 7. Проверка соответствия копии      │   1.1.8.  │
                                     │           │
    дубликату ОСПО                   │   1.1.6.  │  3.1.6.
 ────────────────────────────────────│───────────│─────────
 8. Проверка качества эксплуатацион- │   1.1.9.  │  3.1.7.
                                     │           │
    ных документов                   │           │
 ──────────────────────────────────────────────────────────
 9. Проверка качества маркировки и   │   1.3.1.  │  3.3.1.
                                     │           │
    опломбирования                   │   1.3.2.  │  3.3.2.
                                     │           │
                                     │   1.3.3.  │  3.3.3.
 ────────────────────────────────────│───────────│─────────
 10. Проверка качества упаковки      │   1.4.    │  3.4.
\end{BVerbatim}
\end{center}

\par\textbf{2.2.3.} О готовности ОСПО к испытаниям предприятие-изготовитель должно уведомить представителя заказчика «Извещением», оформленным в установленном на предприятии-изготовителе порядке.\par
К «Извещению» должны быть приложены следующие документы:

Заполненный формуляр на ОСПО;

Комплект программных документов на ОСПО.

\par\textbf{2.2.4.} При обнаружении в процессе приемо-сдаточных испытаний несоответствия ОСПО хотя бы одному требованию настоящих ТУ, этот экземпляр ОСПО считают не выдержавщим испытания и возвращают предприятию-изготовителю для анализа причин дефектов, их устранения и повторной проверки этого экземпляра ОСПО ОТК.\par
При возврате ОСПО представитель заказчика должен составить заключение с изложением причин возврата и ознакомить с ним начальника ОТК.

\par\textbf{2.2.5.} Возвращенные представителем заказчика экземпляры ОСПО могут быть повторно предъявлены к приемо-сдаточным испытаниям после анализа причин дефектов, их устранения и повторной проверки ОТК.\par
\par\textbf{2.2.6.} Приемку экземпляра ОСПО, возвращенного представителем заказчика, допускается возобновлять только при повторном предъявлении «Извещения», на котором должна быть надпись «Вторичное».\par
Вторичное извещение должно быть подписано руководителем предприятия-изготовителя.

К «Извещению» должен прилагаться акт об анализе дефектов, их устранении и повторной проверке ОТК экземпляра ОСПО, возвращенного представителем заказчика, с указанием причин дефектов и мер по их устранению.

Перечень мероприятий по устранению и предупреждению дефектов должен быть согласован с представителем заказчика.

\par\textbf{2.2.7.} Повторные испытания возвращенных представителем заказчика экземпляров ОСПО проводятся в полном объеме, установленном для приемо-сдаточных испытаний.\par
\par\textbf{2.2.8.} Если при повторных испытаниях будет обнаружено несоответствие экземпляра ОСПО требованиям ТУ, то испытания должны быть прекращены, а экземпляр ОСПО возвращен предприятию-изготовителю.\par
Возможность дальнейшего тиражирования и приемки ОСПО решают руководитель предприятия-изготовителя совместно с представителем заказчика и начальником ОТК.

\par\textbf{2.2.9.} На экземпляре ОСПО ( МН или этикетке МН, упаковке ), принятом ОТК и представителем заказчика, должны быть поставлены клейма и пломбы в предусмотренных местах и сделаны соответствующие записи в формуляре ОСПО.\par
\par\textbf{2.2.10.} Окончательно принятым и подлежащим отгрузке считается ОСПО, которое по результатам испытаний соответствует требованиям настоящих ТУ, упаковано, опломбировано ОТК и представителем заказчика и сдано на ответственное хранение предприятию-изготовителю ОСПО.\par
\section*{3.~Методы контроля}\addcontentsline{toc}{section}{3.~Методы контроля}
\subsection*{3.1.~Проверка на соответствие общим требованиям}\addcontentsline{toc}{subsection}{3.1.~Проверка на соответствие общим требованиям}
\par\textbf{3.1.1.} Проверка соответствия изготовленного экземпляра требованиям настоящих ТУ проводится внешним осмотром магнитных носителей.\par
Проверкой должно быть подтверждено, что применены разрешенные виды и типы магнитных носителей.

\par\textbf{3.1.2.} Проверка соответствия изготовленного экземпляра требованиям п.~1.1.3 настоящих ТУ проводится внешним осмотром магнитных носителей и рассмотрением отметки входного контроля магнитных носителей.\par
Проверкой должно быть установлено наличие соответствующих клейм предприятий-изготовителей магнитных носителей, а отметкой входного контроля подтверждается кондиционность магнитных носителей.

\par\textbf{3.1.3.} Для оценки соответствия изготовленного экземпляра требованиям п.~1.1.4 настоящих ТУ устанавливается фактический срок хранения магнитных носителей, который исчисляется от даты изготовления магнитного носителя до даты записи на него информации.\par
Срок хранения магнитных носителей не должен превышать 18 месяцев.

\par\textbf{3.1.4.} Проверка соответствия изготовленного экземпляра п.~1.1.5 настоящих ТУ производится с помощью контрольных проверок ОТК и представителя заказчика технологического процесса изготовления данного экземпляра.\par
Результат последней контрольной проверки должен подтверждать соответствие технологического процесса согласованной и утвержденной инструкции.

\par\textbf{3.1.5.} Проверка изготовленного экземпляра, расположенного на МН, на соответствие требованиям п.~1.1.7 настоящих ТУ производится подсчетом КС по методике, изложенной в приложении.\par
Полученные контрольные суммы должны совпадать с контрольными суммами, которые зафиксированы в формуляре ОСПО.

В случае несовпадения хотя бы одной из контрольных сумм с контрольной суммой, зафиксированной в формуляре ОСПО, должен быть произведен трехкратный подсчет контрольных сумм.

При несовпадении контрольных сумм с контрольными суммами, зафиксированными в формуляре ОСПО ( или при нестабильном получении их значений ) после трехкратного подсчета, считается, что изготовленный экземпляр ОСПО не соответствует требованиям п.~1.1.7 настоящих ТУ и подлежит возврату для анализа причин дефектов, их устранения и перепроверки ОТК.

\par\textbf{3.1.6.} Проверка соответствия изготавливаемого на магнитных носителях экземпляра требованиям п.~1.1.6, п.~1.1.8 настоящих ТУ производится сравнением полученных контрольных сумм с контрольными суммами, которые записаны в формуляре на данный экземпляр.\par
Полученные контрольные суммы должны совпадать с зафиксированными в формуляре.

\par\textbf{3.1.7.} Проверка соответствия программных документов требованиям п.~1.1.9 настоящих ТУ проводится внешним осмотром, предъявленных на испытания эксплуатационных документов ОСПО, выполненных на бумажном носителе.\par
Проверкой должно быть установлено, что после изготовления они не имеют дефектов.

\subsection*{3.2.~Проверка комплектности}\addcontentsline{toc}{subsection}{3.2.~Проверка комплектности}
Проверка соответствия изготовленного экземпляра требованиям п.~1.2 настоящих ТУ производится сравнением предъявленного на испытания комплекта ОСПО с составом, указанным в ТАБЛ.1 ТУ и в формуляре ОСПО.

Комплектность, указанная в ТУ и формуляре, должна иметься в наличии.

\subsection*{3.3.~Проверка маркировки}\addcontentsline{toc}{subsection}{3.3.~Проверка маркировки}
\par\textbf{3.3.1.} Проверка соответствия изготовленного экземпляра, расположенного на МЛ, требованиям п.~1.3.1 настоящих ТУ производится внешним осмотром бобин МЛ.\par
Бобины МЛ должны иметь этикетку установленного на данном предприятии-изготовителе образца.

\par\textbf{3.3.2.} Проверка соответствия ОСПО, расположенного на МЛ, требованиям п.~1.3.2 настоящих ТУ производится внешним осмотром штатной упаковки.\par
Маркировка штатной упаковки должна соответствовать инструкции предприятия-изготовителя, а маркировка транспортной упаковки должна соответствовать технической документации на упаковку завода-изготовителя упаковки.

\par\textbf{3.3.3.} Проверка соответствия изготовленного экземпляра, расположенного на МЛ, требованиям п.~1.1.3. настоящих ТУ производится внешним осмотром МЛ и упаковки.\par
При внешнем осмотре маркировки МЛ и упаковки маркировка должна соответствовать технической документациии на упаковку.

\subsection*{3.4.~Проверка упаковки}\addcontentsline{toc}{subsection}{3.4.~Проверка упаковки}
Проверка соответствия изготовленного экземпляра требованиям п.~1.4 настоящих ТУ производится внешним осмотром его упаковки.

Упаковка должна быть произведена в соответствии с требованиями, изложенными в документации на тару.

\section*{4.~Типовые испытания}\addcontentsline{toc}{section}{4.~Типовые испытания}
\par\textbf{4.1.} Типовые испытания ОСПО проводят в следующих целях:\par
\par\hangindent=1.5em 1) Для оценки эффективности и целесообразности изменений, вносимых в ОСПО или технологию его изготовления, а также для проверки работоспособности ОСПО после внесения изменений, которые затрагивают его функциональные возможности и эксплуатационные характеристики;\par
\par\hangindent=1.5em 2) Для оценки работоспособности ОСПО и определения эксплуатационных характеристик ОСПО после внесения изменений в аппаратуру ВК «Эльбрус 1-КБ», которые затрагивают его функциональные возможности;\par
\par\hangindent=1.5em 3) Для оценки работоспособности и определения эксплуатационных характеристик при создании новых версий ОСПО ВК «Эльбрус 1-КБ».\par
Необходимость проведения типовых испытаний, в зависимости от характера изменений, определяют заказчик или представитель заказчика по согласованию с разработчиком ОСПО.

\par\textbf{4.2.} Типовые испытания могут проводится как на копиях, так и на дубликатах ОСПО, порядок выбора вида и экземпляра ОСПО определяет предприятие-разработчик ОСПО по согласованию с представителем заказчика, что оформляется актом.\par
\par\textbf{4.3.} Типовые испытания ОСПО проводятся по программе и методике, составленной предприятием-разработчиком ОСПО и утвержденной главным конструктором предприятия-разработчика ОСПО и старшим представителем заказчика.\par
\par\textbf{4.4.} Результаты типовых испытаний ОСПО оформляются актом, который утверждается старшим представителем заказчика и главным конструктором предприятия-разработчика ОСПО.\par
\section*{5.~Указания по эксплуатации, транспортированию}\addcontentsline{toc}{section}{5.~Указания по эксплуатации, транспортированию}
\begin{center}И хранению\end{center}
\par\textbf{5.1.} Ввод в эксплуатацию и эксплуатация ОСПО ВК «Эльбрус 1-КБ» должны осуществляться в соответствии с формуляром и эксплуатационной документацией на ОСПО.\par
\par\textbf{5.2.} Поставляемые пользователю МЛ с расположенным на них ОСПО должны использоваться как контрольные экземпляры для изготовления рабочих копий ОСПО, для эксплуатации должны использоваться только рабочие копии ОСПО на МЛ.\par
\par\textbf{5.3.} ОСПО должно эксплуатироваться в нормальных климатических условиях по ГОСТ 20958-80.\par
\par\textbf{5.4.} ОСПО в упаковке может транспортироваться любым видом транспорта на любое расстояние при условии защиты от грязи и атмосферных осадков, при климатических и механических условиях транспортирования в соответствии с ГОСТ 16325-76 и НТД завода-изготовителя (завода-изготовителя тары).\par
\par\textbf{5.5.} Упакованное ОСПО должно храниться в помещении при температуре от 5 градусов до 40 градусов цельсия, относительной влажности воздуха от 30\% до 95\%.\par
\par\textbf{5.6.} Условия транспортирования и хранения изготовленного экземпляра ОСПО на магнитных лентах типа В4502-12 или других типов МЛ с аналогичными характеристиками должны соответствовать требованиям [1].\par
\section*{6.~Гарантии поставщика}\addcontentsline{toc}{section}{6.~Гарантии поставщика}
\par\textbf{6.1.} Предприятие-изготовитель ОСПО гарантирует соответствие изготовленного экземпляра ОСПО требованиям настоящих ТУ при условии его ввода в эксплуатацию предприятием, выполняющим работу по вводу в эксплуатацию и обслуживанию ВК «Эльбрус 1-КБ», а также соблюдения пользователем условий эксплуатации, транспортирования и хранения ОСПО, установленных настоящим ТУ и эксплуатационной документацией.\par
Примечание. В организациях заказчика допускается ввод в эксплуатацию ОСПО силами пользователей ВК «Эльбрус 1-КБ» при условии соблюдения порядка хранения и ввода в эксплуатацию ОСПО согласно эксплуатационной документации на ОСПО.

\par\textbf{6.2.} Предприятие-изготовитель ОСПО гарантирует сохранность и воспроизведение ОСПО ВК «Эльбрус 1-КБ» с магнитных носителей в соответствии с контрольными суммами в течении срока, определяемого сроком действия гарантийных обязательств на ВК.\par
\par\textbf{6.3.} Действие гарантийных обязательств прекращается, если в гарантийный период пользователь внес какие-либо изменения в тексты программ на МЛ, поставляемые заводом-изготовителем ОСПО.\par
\section*{Приложение}\addcontentsline{toc}{section}{Приложение}
\begin{center}Подсчет контрольных сумм магнитных лент\end{center}
1. Контрольная сумма МЛ1, на которой записаны готовые программы операционной системы Диспак, подсчитывается программой переписи с МЛ на МД 2053 при постановке операционной системы. (см. «Формуляр» ИЫ.00110-01 30 01, приложение 1).

По окончании переписи контрольная сумма печатается на АЦПУ.

2. Подсчет контрольных сумм МЛ2, МЛ3 производится при переписи информации с них на магнитный диск 2048 программой «UNICOP» (см. «Формуляр» ИЫ.00110-01 30 01, приложение 1). По окончании переписи на терминал и на АЦПУ выдается контрольная сумма МЛ.

3. Контрольная сумма МЛ4 подсчитывается при переписи с данной МЛ на другую МЛ в режиме ЕС программой «Копес». Программа «Копес» копирует до признака конца ленты (два подряд идущих признака ТМ), поэтому перед копированием необходимо присвоить номера лентам директивами оператора «Есм 1 НУ» --- Ленте чтения и «Есм 2 НУ» --- Ленте записи, здесь НУ --- Номер магнитофона с МЛ.

\begin{center}Структура диалога\end{center}
\begin{center}\small
\begin{BVerbatim}
       -----------------
               Димип 419900
<Версия>       КТ ДОКЭБ 0 ДОКЭБ
=              ред копес
*              ф
*              вый
\end{BVerbatim}
\end{center}

По окончании переписи контрольная сумма МЛ распечатывается на АЦПУ.

\section*{Перечень сокращений}\addcontentsline{toc}{section}{Перечень сокращений}
АЦПУ --- Алфавитно-цифровое печатающее устройство

ВК --- Вычислительный комплекс

КС --- Контрольная сумма

МД --- Магнитный диск

МЛ --- Магнитная лента

МН --- Магнитный носитель

ОСПО --- Общее системное программное обеспечение

ОТК --- Отдел технического контроля

ТМ --- ТАРЕ-MARK (МГЗ) ленточный маркер

ТУ --- Технические условия

Эд --- Эксплуатационная документация

\section*{Перечень ссылочных документов}\addcontentsline{toc}{section}{Перечень ссылочных документов}
\section*{1.~ТУ6-17-686-75 лента магнитная тип В4502-12.ТЕХНИЧЕСКИЕ}\addcontentsline{toc}{section}{1.~ТУ6-17-686-75 лента магнитная тип В4502-12.ТЕХНИЧЕСКИЕ}
\begin{center}Условия.\end{center}
\section*{2.~ГОСТ 20958-80 машины вычислительные.ленты магнитные}\addcontentsline{toc}{section}{2.~ГОСТ 20958-80 машины вычислительные.ленты магнитные}
\begin{center}Без записи. Общие технические условия.\end{center}
\section*{3.~ГОСТ 16325-76 машины вычислительные электронные циф-}\addcontentsline{toc}{section}{3.~ГОСТ 16325-76 машины вычислительные электронные циф-}
\begin{center}Ровые общего назначения.общие техни-\end{center}
\begin{center}Ческие требования.\end{center}
\section*{4.~ГОСТ 19.106-78 единая система программной документации}\addcontentsline{toc}{section}{4.~ГОСТ 19.106-78 единая система программной документации}
\begin{center}Требования к программным документам,\end{center}
\begin{center}Выполненным печатным способом.\end{center}
\section*{5.~Вычислительный комплекс «Эльбрус 1-КБ»}\addcontentsline{toc}{section}{5.~Вычислительный комплекс «Эльбрус 1-КБ»}
Общее системное программное обеспечение.

Формуляр. ИЫ.00110-01 30 01

\section*{Лист регистрации изменений}\addcontentsline{toc}{section}{Лист регистрации изменений}
.-----------------------------------------------------------------.

\begin{center}\small\begin{tabularx}{\linewidth}{|>{\raggedright\arraybackslash\hsize=0.176\hsize}X|>{\raggedright\arraybackslash\hsize=0.706\hsize}X|>{\raggedright\arraybackslash\hsize=1.059\hsize}X|>{\raggedright\arraybackslash\hsize=1.059\hsize}X|>{\raggedright\arraybackslash\hsize=1.059\hsize}X|>{\raggedright\arraybackslash\hsize=1.059\hsize}X|>{\raggedright\arraybackslash\hsize=1.235\hsize}X|>{\raggedright\arraybackslash\hsize=1.588\hsize}X|>{\raggedright\arraybackslash\hsize=1.941\hsize}X|>{\raggedright\arraybackslash\hsize=1.235\hsize}X|>{\raggedright\arraybackslash\hsize=0.706\hsize}X|>{\raggedright\arraybackslash\hsize=0.176\hsize}X|}\hline
 & \small\bfseries\hspace{0pt}И з м. & \small\bfseries\hspace{0pt}Номер измененных & \small\bfseries\hspace{0pt}Лист замененных & \small\bfseries\hspace{0pt}В (СТ новых & \small\bfseries\hspace{0pt}Аниц) аннулированных & \small\bfseries\hspace{0pt}Всего листов (СТР.) В документе & \small\bfseries\hspace{0pt}Номер документа & \small\bfseries\hspace{0pt}Входящий номер сопроводи-тельного документа & \small\bfseries\hspace{0pt}Подпись & \small\bfseries\hspace{0pt}Д а т а &  \\ \hline
\end{tabularx}\end{center}

\end{document}
